% Biografias copiadas do modelo ANA-ML 2.0 (access (2).tex); as fotos devem estar nesta pasta.
\begin{IEEEbiography}[{\includegraphics[width=1in,height=1.5in,clip,keepaspectratio]{marx_viana_pb.png}}]{Marx Viana} received a B.S. in computer science from Federal University of Viçosa (UFV), in 2010, and the Master's in Informatics and Ph.D. degrees at the Pontifical Catholic University of Rio de Janeiro (PUC-RJ), in 2012 and 2016, respectively.

He is a currently a researcher at PUC-RJ. His research in computer science span a wide range of topics with an emphasis on software engineering, and focus on software frameworks, object-oriented design and programming, software agents, software adaptation, and normative systems.
\end{IEEEbiography}

\begin{IEEEbiography}[{\includegraphics[width=1in,height=1.25in,clip,keepaspectratio]{elder-menino.jpg}}]{Elder Cirilo} received a B.S. in Computer Science from Federal University of Juiz de Fora (UFJF), in 2006, and the Master's and Ph.D. degrees in Informatics at the Pontifical Catholic University of Rio de Janeiro (PUC-RJ), in 2008 and 2012, respectively. He has been an Assistant Professor in the Department of Computer Science at the Federal University of São João del Rei since 2013.

He has experience in the area of software engineering, with an emphasis on the following themes: configurable software systems, domain specific languages, feature-oriented software development, end-user programming and experimental software engineering.
\end{IEEEbiography}

\begin{IEEEbiography}[{\includegraphics[width=1in,height=1.25in,clip,keepaspectratio]{palencar.png}}]{Paulo Alencar} is a Research Professor with the David R. Cheriton School of Computer Science at the University of Waterloo, and Associate Director of the Computer Systems Group (CSG).

Dr. Alencar has received international research awards from organizations such as Compaq and IBM. He has published over 250 refereed publications and has been a member of program committees of numerous highly-regarded conferences and workshops. His recent research in information technology and software engineering has focused on high-level software architectures, design, components and their interfaces; software frameworks and application families; software processes, automated workflows and work graphs, and evolution; web-based approaches and applications; open and big data applications; context-aware and event-based systems; software agents;
machine learning; cognitive chatbots; artificial intelligence; and formal
methods. He has been Principal or Co-Principal Investigator in projects
supported by NSERC, ORF-RE, IBM, SAP, CITO, CSER, Bell, and funding agencies in Canada U.S., Brazil, Germany and Argentina. He is a member of the Association of Computing Machinery (ACM), Institute of Electrical and Electronic Engineers (IEEE), Association for the Advancement of Artificial Intelligence (AAAI), Waterloo Water Institute (part of the Global Water Futures initiative), and Waterloo Artificial Intelligence Institute.
\end{IEEEbiography}

\begin{IEEEbiography}[{\includegraphics[width=1in,height=1.25in,clip,keepaspectratio]{don_cowan.jpeg}}]{Don Cowan} (Life Member, IEEE) is currently a Distinguished Professor Emeritus of computer science with the University of Waterloo and the Director of the Computer Systems Group. He has made contributions to computer science in areas, such as computer science, software engineering, and complex applications. He is the author and coauthor or an editor of over 300 refereed articles and 17 books in computer/communications, software engineering, education, environmental information systems, and mathematics. He has supervised over 120 graduate students and postdoctoral fellows. His group has developed over 80 Web-based and mobile software systems for many application areas. The University of Waterloo recognized his contributions to development of graduate students by presenting him with the Award of Excellence in Graduate Supervision. He was recognized for his research and support of the development of computer science in Brazil by being awarded the National Order of Scientific Merit (Grand Cross), the country's highest civilian scientific honour, by the President of Brazil. In 2009, he received the Waterloo Award, the City of Waterloo's highest civic honour, for his contributions to the City of Waterloo. In 2010, he was named a Distinguished Scientist by the Association for Computing Machinery. In 2011, he received a Doctor of Science (Honoris Causa) from the University of Guelph for his contributions to Computer Science and Software Engineering.
\end{IEEEbiography}

\begin{IEEEbiography}[{\includegraphics[width=1in,height=1.25in,clip,keepaspectratio]{carlos_lucena.jpeg}}]{Carlos Lucena} received a BSc degree from the Pontifical Catholic University of Rio de Janeiro (PUC-Rio), Brazil, in 1965,
a Math degree in computer science from the University of Waterloo, Canada, in 1969, and a PhD degree in computer science from the University of California at Los Angeles in 1974. He has been a full professor in the Department of Informatics at PUC-Rio
since 1982.

Dr. Carlos Lucena has made extensive and internationally recognized research contributions, being the author or co-author of over 600 refereed papers and 19 books in the area of software design and formal methods in software
engineering, his primary areas of research. He has also been a member of the program committees of more than
62 national and international conferences as well as member of the editorial board of some important journals in his area of research. Until 2015, he has supervised 45 PhD theses and 114 M.Sc. dissertations. His former PhD students are faculty members in universities in Brazil and abroad (e.g., USA and Canada). He is also a member of
the Editorial Board of the International Journal on Formal Aspects of Computing (New York: Springer-Verlag).
He worked at the Pontifical Catholic University of Rio de Janeiro as vice-rector of the University, Dean for the Center for Science and Technology, and several times as Director of the Computer Science Department. Prof.
Lucena was awarded the insignia for the Grand Cross of the Order of Scientific Merit of the Presidency of the
Republic of Brazil, the Medal of Scientific Merit Carlos Chagas Filho from the Council of Directors of FAPERJ (State Research Council), the Alvaro Alberto Award for Science and Technology (Ministry of Science and Technology)
and several times the IBM Innovation Award, among many others. Prof. Lucena is a Fellow of the Guggenheim
Foundation, an ACM Fellow and a member of the Brazilian Academy of Sciences.
\end{IEEEbiography}
